\documentclass[10pt,a4paper]{amsart}
\usepackage[margin=19mm]{geometry}
\usepackage{amsmath,amssymb,mathtools}
\usepackage[scr=rsfs,cal=euler]{mathalfa}
\usepackage{xfrac}
\usepackage[unicode,hidelinks,bookmarks=false]{hyperref}
\newcommand{\ie}{i.e.\ }
\newcommand{\cf}{cf.\ }
\newcommand{\resp}{respectively\ }
\newcommand{\Subsection}{\subsection}
\providecommand{\nobreakdash}{\nobreak\hbox{-}}
\setlength{\emergencystretch}{2em}
\multlinegap0pt
\usepackage{bm}
\usepackage{bbm}
\usepackage{dsfont}
\usepackage{xspace}
\usepackage{cancel}
\usepackage{mathtools}
\usepackage{cleveref}
\usepackage{amsthm}
\makeatletter
\@ifundefined{theorem}{\newtheorem{theorem}{Theorem}}{}
\@ifundefined{lemma}{\newtheorem{lemma}[theorem]{Lemma}}{}
\@ifundefined{observation}{\newtheorem{observation}[theorem]{Observation}}{}
\makeatother
\newcommand{\E}{\mathbb{E}}
\newcommand{\F}{\mathcal{F}}
\renewcommand{\Pr}{\mathbb{P}}
\title[Frequent elements in union-closed set families]{Frequent elements in union-closed set families}
\author{Shagnik Das, Saintan Wu}
\date{Source: arXiv:2412.03862v3 (CC BY 4.0)}
\begin{document}
\maketitle

\noindent\emph{Source and licence.} Frequent elements in union-closed set families. Version arXiv:2412.03862v3 (CC BY 4.0), distributed under the Creative Commons Attribution 4.0 International licence, as recorded in the arXiv licence metadata for this exact version and shown on the abstract page https://arxiv.org/abs/2412.03862v3. Full rights notice: licenses/arxiv-2412.03862-CC-BY-4.0.txt.\par\smallskip

\noindent\textit{Editorial scope.} Counted unit: the Theorem whose source label is \texttt{thm:largeNew} in the article (source0.tex lines 262--265, source label \texttt{thm:largeNew}), delivered together with the article's own complete original proof of it (source0.tex lines 266--298, one \texttt{proof} environment of 33 lines). No mathematics is written, repaired or re-derived: statement and proof are the article's own text line for line. Required context retained uncounted: obs:equivalence (lines 151--153); lemma:afterk (lines 251--255). Every definition, display and label of those lines is retained unchanged. Omitted from this package and not redistributed: 1--150, 154--250, 256--261, 299--424. Nothing inside a retained excerpt is omitted or altered: every retained body is delivered exactly as the article's lines are written, so no omission receipt of any kind accompanies this package. Citations: the retained excerpts carry no citation command at all, so no bibliography entry of the article is transcribed and no reference is redistributed. Reference adaptation: none was needed and none was made; every reference inside the delivered unit resolves inside it, so no unresolved reference or citation is produced. Printed slips kept literal: no printed word, symbol, hypothesis or number of the counted statement or of its proof was altered. Layout and class: the article's own class is \texttt{article} with packages including babel, geometry, type1cm, titlesec, fancyhdr, xcolor, graphicx, titling, tabularx, enumitem, amsmath, amsthm, amssymb, mathtools; the wrapper is the lane's pinned \texttt{amsart} editorial wrapper at 10pt on A4, which restates the theorem-like environments the retained text needs and adds the article's own macro definitions that the retained text needs (\texttt{\textbackslash E}, \texttt{\textbackslash F}, \texttt{\textbackslash Pr}), with extra wrapper packages (bm, bbm, dsfont, xspace, cancel, mathtools, cleveref). The delivered numbering is the unit's own. Assets: no retained span contains a figure environment, a graphics-inclusion command, a PDF-page inclusion, a table or a TikZ picture; no image file is copied into this package, the package declares no asset dependency and no image file is redistributed with it. Licence: version arXiv:2412.03862v3 of this article is distributed under the Creative Commons Attribution 4.0 International licence, as recorded in the arXiv licence metadata for this exact version and shown on the abstract page https://arxiv.org/abs/2412.03862v3. A full rights notice is delivered at licenses/arxiv-2412.03862-CC-BY-4.0.txt. Characters: every retained excerpt is pure ASCII, so the delivered engine, XeLaTeX with its default Latin Modern text font, reports no missing character.
\begin{observation} \label{obs:equivalence}
    Assuming the Union-Closed Sets Conjecture, Nagel's conjecture is true for all $k$.
\end{observation}

    \begin{lemma}\label{lemma:afterk}
        Let $A,B$ be i.i.d. random sets on the family $2^{[n]}$. Suppose further that $\E[A_i]\leq \alpha< \frac{3-\sqrt{5}}{2}$. Then 
        \[H((A\cup B)_i|A_{< i},B_{<i}) \geq \lambda_\alpha H(A_i |A_{< i}),\]
        where $\lambda_\alpha =\frac{H(2\alpha-\alpha^2)}{H(\alpha)}>1$.\footnote{Sawin also showed that this inequality holds when $\alpha \ge \frac{3 - \sqrt{5}}{2}$, but with the constant $\lambda_{\alpha}$ defined to be $\frac{1+\sqrt{5}}{2} (1 - \alpha)$ instead of $\frac{H(2\alpha - \alpha^2)}{H(\alpha)}$.} In particular, $H(A\cup B)\geq \lambda_\alpha H(A)$.
    \end{lemma}

    \begin{theorem}\label{thm:largeNew}
        Let $k\geq 2$ and $0 < \alpha<\frac{3-\sqrt{5}}{2}$. If $\F$ is a union-closed family with $f_k(\F)\leq \alpha$, then $\log_2|\F|\leq \frac{\lambda_\alpha}{\lambda_\alpha-1}(k-1)$. In fact,
        \[\log_2|\F| \leq \frac{\lambda_\alpha}{\lambda_\alpha-1}\cdot \frac{2^{k-1}(k-1)}{2^{k-1}-1} -\frac{1}{(\lambda_\alpha-1)}\log_2\left(\frac{\lambda_\alpha2^{k-1}(k-1)e}{(2^{k-1}-1)\log_2 e }\right).\]
    \end{theorem}

	\begin{proof}
            Let $\F$ be a union-closed family with $f_k(\F) \le \alpha$. As we do not have information about the frequencies of the $k-1$ most frequent elements, we shall project onto their complement, using the map $\pi_{k-1}$ from the proof of~\cref{obs:equivalence}:
	\begin{align*}
		\pi_{k-1}: 2^{[n]}&\to 2^{[n]\setminus[k-1]}\\
		F&\mapsto F\setminus [k-1].
	\end{align*}
	
        Now let $X$ be a uniformly random set in $\F$, and let $A = \pi_{k-1}(X)$ be its projection. Observe that for any $i \ge k$, we have $\E[A_i] = \Pr(i \in A) = \Pr(i \in X) \le \alpha$. Hence, if we let $B$ be independent of and identically distributed as $A$,~\cref{lemma:afterk} gives $H(A\cup B)\geq \lambda_\alpha H(A)$, where $\lambda_{\alpha}=\frac{H(2\alpha-\alpha^2)}{H(\alpha)}$. To obtain the desired contradiction, we need provide an upper bound for $H(A\cup B)$ and a lower bound for $H(A)$. 	
	
	For the upper bound, we simply apply the support bound. Since $\pi_{k-1}(\F)$ is still union-closed, $A\cup B$ is again a distribution on $\pi_{k-1}(\F)$. Hence, $H(A\cup B)\leq \log_2|\pi_{k-1}(\F)|$. 
	
	For a lower bound on $H(A)$, observe that, since $X$ is uniformly distributed over $\F$, we have $H(X) = \log_2 |\F|$. On the other hand, we have $H(X)=H(A,X)$, since $A$ is determined by $X$. Applying the chain rule gives $H(A,X)=H(A)+H(X|A)$. Putting this all together yields $H(A)=\log_2|\F|-H(X|A)$.

    Combining the upper and lower bounds results in
    \begin{equation} \label{eqn:entropy}
        \log_2 |\pi_{k-1}(\F)| \ge \lambda_\alpha \left( \log_2 |\F| - H(X|A) \right).
    \end{equation}
    
    We next evaluate $H(X|A)$. By definition, 
    \[ H(X|A) = \sum_{F \in \pi_{k-1}(\F)} \Pr(A = F)H(X|A = F) = \sum_{F \in \pi_{k-1}(\F)} \frac{|\pi_{k-1}^{-1}(F)|}{|\F|} H(X|A=F). \] 
    Conditioning on $A = F$, we have that $X$ is uniformly distributed on $\pi_{k-1}^{-1}(F)$, and so 
    \[ H(X|A) = \sum_{F \in \pi_{k-1}(\F)} \frac{|\pi_{k-1}^{-1}(F)|}{|\F|} \log_2 |\pi_{k-1}^{-1}(F)|. \]

    For a simple bound, observe that $|\pi_{k-1}^{-1}(F)| \le 2^{k-1}$ for all $F$, and so it follows that $H(X|A) \le k-1$. Substituting this into~\eqref{eqn:entropy}, together with the trivial bound $|\pi_{k-1}(\F)| \le |\F|$, we get $\log_2 |\F| \le \frac{\lambda_{\alpha}}{\lambda_{\alpha} - 1} (k-1)$, proving the first inequality from the statement of the theorem.

    For the more precise estimate, note that the function $f(x) = x \log_2 x$ is convex, and hence for any $x \in [1, 2^{k-1}]$, we have $f(x) \le \frac{2^{k-1}-x}{2^{k-1} - 1} f(1) + \frac{x - 1}{2^{k-1}-1} f(2^{k-1}) = (x-1)\frac{2^{k-1}(k-1)}{2^{k-1}-1}$. Thus, $H(X|A) = \sum_{F \in \pi_{k-1}(\F)} \frac{f(|\pi_{k-1}^{-1}(F)|)}{|\F|} \le \sum_{F \in \pi_{k-1}(\F)} \frac{|\pi_{k-1}^{-1}(F)|-1}{|\F|} \cdot \frac{2^{k-1}(k-1)}{2^{k-1} - 1} = \left( 1 - \frac{|\pi_{k-1}(\F)|}{|\F|} \right) \frac{2^{k-1}(k-1)}{2^{k-1}-1}.$

    Letting $\rho = \frac{|\pi_{k-1}(\F)|}{|\F|}$, we can substitute this into~\eqref{eqn:entropy} to obtain
    \[ \log_2 \rho + \log_2 |\F| \ge \lambda_\alpha \left( \log_2 |\F| - (1 - \rho) \frac{2^{k-1}(k-1)}{2^{k-1} - 1} \right), \]
    which can be rearranged to give $(\lambda_\alpha - 1) \log_2 |\F| \le \log_2 \rho + (1 - \rho) \lambda_\alpha \frac{2^{k-1}(k-1)}{2^{k-1}-1}$.
    
    By differentiating, we find the right-hand side is maximised when $\rho = \frac{(2^{k-1}-1) \log_2 e}{\lambda_\alpha 2^{k-1} (k-1)}$, for which the right-hand side becomes $\frac{\lambda_\alpha 2^{k-1}(k-1)}{2^{k-1}-1} - \log_2 \left( \frac{\lambda_\alpha 2^{k-1}(k-1)e}{(2^{k-1}-1)\log_2 e}\right)$. This gives the desired upper bound on $\log_2 |\F|$.
    \end{proof}

\end{document}
